\subsection{Resultados de la señal 1}
\label{sec:resultados}

\subsubsection{Señal recibida y compresión}
La \autoref{fig:recibida} muestra un momento de muestreo (el pulso $t_0$). Hasta el rango mínimo de 1,5 km la señal es nula, porque el receptor está bloqueado durante la transmisión y la recuperación. Entre 1,5 y 2,5 km aparece un eco intenso, de hasta $6{,}9 \times 10^{-6}$, que corresponde a clutter cercano. Más allá, los ecos de los blancos quedan por debajo del ruido: un chirp de 10 $\mu$s extiende cada eco sobre 1,5 km y su amplitud es comparable a la del ruido.

\begin{figure}[H]
    \centering
    \includegraphics[width=0.85\textwidth]{\figdir/01_senal_recibida.png}
    \caption{Señal 1: un momento de muestreo de la señal recibida. Arriba, partes real e imaginaria; abajo, módulo, en función del rango.}
    \label{fig:recibida}
\end{figure}

Después del filtro apareado (\autoref{fig:compresion}), la energía de cada eco se concentra en un pico de unos 75 m y emergen claramente ocho blancos, entre 10 y 22,2 km, con amplitudes de 4 a 15 veces el nivel del ruido. Esto ilustra la ganancia de compresión ($BT = 20$, 13 dB) y la resolución de 75 m: los blancos de 15,21 y 15,51 km, separados 300 m, aparecen como dos picos distintos.

\begin{figure}[H]
    \centering
    \includegraphics[width=0.85\textwidth]{\figdir/02_compresion.png}
    \caption{Señal 1: módulo de la señal sin comprimir (arriba) y comprimida con el filtro apareado (abajo) para el mismo momento de muestreo.}
    \label{fig:compresion}
\end{figure}

\subsubsection{Canceladores y CFAR}
Las figuras \ref{fig:cfar_mti} a \ref{fig:cfar_mtidc} muestran, para cada cancelador:
\begin{itemize}
    \item las señales sin comprimir de los momentos de muestreo que usa;
    \item las señales comprimidas;
    \item la salida del cancelador con el umbral CFAR (ganancia 3,93);
    \item los cruces por cero que marcan el inicio ($+2$) y el fin ($-2$) de cada detección.
\end{itemize}

\begin{figure}[H]
    \centering
    \includegraphics[width=0.88\textwidth]{\figdir/04_cfar_mti_sc.png}
    \caption{Señal 1: CFAR del MTI simple cancelador ($x_1 - x_0$).}
    \label{fig:cfar_mti}
\end{figure}

\begin{figure}[H]
    \centering
    \includegraphics[width=0.88\textwidth]{\figdir/05_cfar_sti_sc.png}
    \caption{Señal 1: CFAR del STI simple cancelador ($x_1 + x_0$).}
    \label{fig:cfar_sti}
\end{figure}

\begin{figure}[H]
    \centering
    \includegraphics[width=0.88\textwidth]{\figdir/06_cfar_mti_dc.png}
    \caption{Señal 1: CFAR del MTI doble cancelador ($x_2 - 2x_1 + x_0$).}
    \label{fig:cfar_mtidc}
\end{figure}

\begin{table}[H]
    \centering
    \footnotesize
    \caption{Señal 1: rangos detectados por cada cancelador [km].}
    \label{tab:det_canceladores}
    \begin{tabular}{|l|p{6.2cm}|p{5.2cm}|}
        \hline
        \textbf{Cancelador} & \textbf{Un momento de muestreo ($t_0$, $t_1$, $t_2$)} & \textbf{Integrado sobre todos los pulsos} \\
        \hline
        MTI simple & 10,02 -- 15,21 -- 20,01 -- 22,23 & 10,02 -- 15,21 -- 20,01 -- 22,23 \\
        \hline
        STI simple & 12,02 -- 15,51 -- 17,61 -- 20,01 -- 21,02 -- 22,23 & 12,02 -- 17,61 -- 20,01 -- 21,02 -- 22,23 \\
        \hline
        MTI doble & 10,02 -- 15,21 -- 20,01 -- 22,25 & 10,02 -- 15,21 -- 20,01 -- 22,23 \\
        \hline
    \end{tabular}
\end{table}

\textbf{Análisis:}
\begin{itemize}
    \item \underline{MTI simple y doble:} detectan los mismos cuatro blancos móviles (10,02, 15,21, 20,01 y 22,23 km). Cancelan el clutter cercano y los blancos fijos de 12,02, 17,61 y 21,02 km, que en la señal comprimida son de los ecos más intensos. El doble cancelador da prácticamente el mismo resultado, con un nulo más ancho alrededor de la velocidad cero.
    \item \underline{STI simple:} refuerza los ecos fijos y detecta los tres blancos fijos. No detecta el de 10,02 km: su frecuencia Doppler ($-742$ Hz) está a menos de 8 Hz de $-PRF/2$, donde el STI tiene su nulo (atenuación de unos 30 dB, \autoref{fig:canceladores}).
    \item \underline{Enmascaramiento:} el blanco de 15,51 km no aparece en el MTI. Está a solo 300 m del de 15,21 km, que es más intenso en el MTI y eleva el umbral CFAR a su alrededor, como se ve en la \autoref{fig:cfar_mti}. Tampoco aparece en el STI integrado. El procesamiento Doppler lo separa, porque tiene otra velocidad.
\end{itemize}

\subsubsection{Procesamiento Doppler}
La \autoref{fig:doppler3d} es el gráfico 3D del producto Doppler de la señal MTI: cada pico es un blanco móvil ubicado en su rango y su velocidad. La \autoref{fig:mapa} muestra la vista superior en dB con los blancos detectados.

\begin{figure}[H]
    \centering
    \includegraphics[width=0.85\textwidth]{\figdir/07_doppler_3d.png}
    \caption{Señal 1: procesamiento Doppler de la señal del MTI simple cancelador (módulo en función de la velocidad y el rango).}
    \label{fig:doppler3d}
\end{figure}

\begin{figure}[H]
    \centering
    \includegraphics[width=0.85\textwidth]{\figdir/08_mapa_rango_velocidad.png}
    \caption{Señal 1: mapa rango-velocidad en dB con los blancos móviles (cruces) y los blancos fijos, ubicados en velocidad nula (cuadrados).}
    \label{fig:mapa}
\end{figure}

\subsubsection{Blancos detectados}
La \autoref{tab:blancos} reúne el resultado final: \textbf{ocho blancos, cinco móviles y tres fijos}.

\begin{table}[H]
    \centering
    \footnotesize
    \caption{Señal 1: blancos detectados. La velocidad positiva indica acercamiento y la negativa, alejamiento. Las ganancias del MTI y el STI corresponden a la frecuencia Doppler de cada blanco.}
    \label{tab:blancos}
    \begin{tabular}{|c|r|r|r|r|r|r|r|l|}
        \hline
        \textbf{N.°} & \textbf{Rango [km]} & \textbf{$v_r$ [m/s]} & \textbf{$v_r$ [km/h]} & \textbf{$f_D$ [Hz]} & \textbf{SNR [dB]} & \textbf{MTI [dB]} & \textbf{STI [dB]} & \textbf{Tipo} \\
        \hline
        1 & 10,02 & $-85{,}7$ & $-308$ & $-742$ & 32,8 & $+6{,}0$ & $-29{,}9$ & Móvil (*) \\
        2 & 12,02 & 0 & 0 & 0 & -- & $-\infty$ & $+6{,}0$ & Fijo \\
        3 & 15,21 & $+40{,}2$ & $+145$ & $+349$ & 38,9 & $+2{,}5$ & $+3{,}5$ & Móvil \\
        4 & 15,51 & $-22{,}7$ & $-82$ & $-197$ & 37,9 & $-1{,}9$ & $+5{,}3$ & Móvil \\
        5 & 17,61 & 0 & 0 & 0 & -- & $-\infty$ & $+6{,}0$ & Fijo \\
        6 & 20,01 & $-43{,}7$ & $-157$ & $-379$ & 42,2 & $+3{,}1$ & $+2{,}9$ & Móvil \\
        7 & 21,02 & 0 & 0 & 0 & -- & $-\infty$ & $+6{,}0$ & Fijo \\
        8 & 22,23 & $-31{,}5$ & $-113$ & $-273$ & 39,2 & $+0{,}7$ & $+4{,}5$ & Móvil \\
        \hline
    \end{tabular}

    \vspace{0.1cm}
    \raggedright \scriptsize (*) Velocidad en el borde del intervalo no ambiguo ($\pm 86{,}5$ m/s): podría ser una velocidad replegada, por ejemplo $+87{,}4$ m/s, que difiere en $\lambda \, PRF/2 = 173$ m/s.
\end{table}

\begin{itemize}
    \item \underline{Blancos fijos:} los de 12,02, 17,61 y 21,02 km tienen una salida del STI entre 24 y 27 dB mayor que la del MTI. Son ecos sin Doppler, como blancos detenidos o reflectores fijos, que el MTI cancela.
    \item \underline{Residuos y falsas alarmas:} el mapa rango-Doppler tiene detecciones débiles en velocidad nula a 20,01 y 22,23 km. Son unos 25 dB menores que los blancos de esas celdas y corresponden a la dispersión Doppler por la fluctuación de su RCS. Además hay tres detecciones aisladas (1,5, 6,2 y 13,7 km) de unos 10 dB de SNR que ningún cancelador confirma. Todas se descartaron con los criterios de la \autoref{sec:operaciones}.
    \item \underline{Resolución en velocidad:} es de 1,75 m/s. Las velocidades informadas tienen una incertidumbre de media celda Doppler ($\pm 0{,}9$ m/s).
\end{itemize}
